Two Romanian astronomers Ovidiu Tercu and Alex Dumitriu from The Astronomical
Observatory of the Museum Complex of Natural Sciences in Galati, Romania,
recently discovered --- in September 2013 --- two variable stars. They used
for that a telescope with the main mirror diameter of 40\,cm and a SBIG
STL-6303e CCD camera.

With the accord of the AAVSO (American Association of Variable Stars
Observers), the two stars have now Romanian names: Galati V1 and respectively
Galati V2. The two stars are circumpolar, located in Cassiopeia and
respectively in Andromeda constellation. The two stars are visible above the
horizon, from the territory of Romania, all over the year. The galactic
coordinates of the two stars are: Galati V1 ($G_1 = 114.371^\circ$;
$g_1 = -11.35^\circ$) and Galati V2 ($G_2 = 113.266^\circ$;
$g_2 = -16.177^\circ$).

Another star, discovered by the Romanian astronomer Nicolas Sanduleak, has
also a Romanian name --- Sanduleak $-69^\circ$ 202; it explodes as the
supernova SN 1987. This star was localized in the Dorado constellation from
the Large Magellan Cloud, by the coordinates:
\[
\alpha = 5^{\mathrm{h}}35^{\mathrm{min}}28.03^{\mathrm{s}};\quad
\delta = -69^\circ 16' 11.79'';\quad
G = 279.7^\circ;\quad
g = -31.9^\circ.
\]

Estimate the angular distance between the stars Galati V1 and Galati V2.
